% Gerado por scripts/etapa1_sazonalidade_anual.py — requer \usepackage{booktabs}
\begin{table}[htbp]
\centering
\small
\caption{Índice sazonal mensal}
\label{tab:saz_indice_mensal}
\begin{tabular}{lrrrrrr}
\toprule
Mês & 2011 & 2012 & 2013 & 2014 & 2015 & Média dos anos completos \\
\midrule
Jan & -- & 1,000 & 0,943 & 0,940 & 0,952 & 0,961 \\
Fev & -- & 1,029 & 0,954 & 0,947 & 0,934 & 0,977 \\
Mar & -- & 1,024 & 0,964 & 0,980 & 0,951 & 0,989 \\
Abr & -- & 1,006 & 0,940 & 1,039 & 0,949 & 0,995 \\
Mai & -- & 1,019 & 0,963 & 0,993 & 0,957 & 0,992 \\
Jun & -- & 1,071 & 1,075 & 1,043 & 1,039 & 1,063 \\
Jul & -- & 1,036 & 1,076 & 1,049 & 1,052 & 1,054 \\
Ago & 0,934 & 1,085 & 1,114 & 1,068 & 1,075 & 1,089 \\
Set & 0,973 & 1,048 & 1,105 & 1,038 & 1,053 & 1,064 \\
Out & 1,056 & 0,949 & 0,991 & 0,982 & -- & 0,974 \\
Nov & 0,989 & 0,897 & 0,923 & 0,979 & -- & 0,933 \\
Dez & 1,060 & 0,971 & 0,950 & 0,924 & -- & 0,948 \\
\bottomrule
\end{tabular}
\par\smallskip
{\footnotesize Nota: FOODS, lojas da Califórnia, dias disponíveis e de loja aberta; demanda média por série-dia. Índice = demanda do mês ÷ média móvel centrada 2×12 meses (remove tendência e mudanças de nível). 1,05 = 5\% acima da tendência. Só meses completos; anos nas pontas ficam parciais.}
\end{table}
